El efecto Skin es un fenómeno físico que se manifiesta al circular corriente alterna por un conductor, de manera que la corriente se concentra en su superficie, disminuyendo hacía el interior conforme aumenta la frecuencia de la señal. A partir de las ecuaciones de Maxwell en un medio conductor homogéneo y sin densidad de carga libre es posible explicar el fenómeno, donde la solución a la ecuación de onda es modificada con un término exponencial que depende de factores como el número de onda, la distancia en el medio $k$, la  conductividad $\sigma$, la permeabilidad magnetica $\mu$ y la frecuencia angular $\omega$. En materiales donde la conductividad es mayor que $\varepsilon \omega$, siendo $\varepsilon$ la permeabilidad electrica, el término exponencial se reduce a \cite{ecuacionesSkinEffect} 
\begin{equation}
    \psi(z) = \exp\left[ -z/(2/\mu \sigma \omega)^{1/2} \right]
    \label{eq: psi(z)}
\end{equation}
definiendo al termino $\delta = \sqrt{\frac{2}{\mu \sigma \omega}}$ como la profundidad de penetración de la corriente. En este experimento se estudiará el campo inducido entre dos bobinas separadas por una placa metálica conductora con el objetivo de analizar la atenuación producida por el efecto skin en función de la frecuencia y el espesor de la placa, obteniendo así la profundidad de penetración de la corriente y la conductividad del material.